% Generated draft: requires mathematical and source review.
\documentclass[11pt]{article}
\usepackage[margin=25mm]{geometry}
\usepackage{fontspec}
\setmainfont{Latin Modern Roman}
\usepackage{amsmath,amssymb,mathtools,unicode-math}
\setmathfont{Latin Modern Math}
\usepackage{enumitem,needspace,xurl,hyperref}
\setlength{\parindent}{0pt}
\setlength{\parskip}{5pt}
\setlength{\emergencystretch}{3em}
\newcommand{\N}{\mathbb{N}}
\newcommand{\Z}{\mathbb{Z}}
\newcommand{\Q}{\mathbb{Q}}
\newcommand{\R}{\mathbb{R}}
\newcommand{\C}{\mathbb{C}}
\newcommand{\F}{\mathbb{F}}
\newcommand{\A}{\mathbb{A}}
\newcommand{\PP}{\mathbb{P}}
\newcommand{\E}{\mathbb{E}}
\newcommand{\eps}{\varepsilon}
\newcommand{\dd}{\,\mathrm{d}}
\begin{document}
\section*{UnsolvedMath ID: 701079}
\textbf{Source problem code:} catalogue entry
\subsection{Definitions and mathematical statement}
Let \(p\) be a prime and let \(P\) be a finite \(p\)-group. Define its lower central series by \(\gamma_1(P)=P\) and \(\gamma_{i+1}(P)=[\gamma_i(P),P]\). Its nilpotency class \(c(P)\) is the least integer \(c\geq 0\) such that \(\gamma_{c+1}(P)=1\); thus the trivial group has class \(0\). For \(x\in P\), let \(x^P=\{g^{-1}xg:g\in P\}\). Define the breadth \(b(P)\) by
\[
p^{b(P)}=\max_{x\in P}|x^P|=\max_{x\in P}[P:C_P(x)],
\]
where \(C_P(x)=\{g\in P:gx=xg\}\).

\textbf{Class--Breadth Problem (catalogue entry).} Is it true that, for every odd prime \(p\) and every finite \(p\)-group \(P\),
\[
c(P)\leq b(P)+1?
\]

\textbf{Supplied contextual assertions.} The record describes
\[
c(P)<\frac{p}{p-1}b(P)+1
\]
as the best known bound and attributes it to C. R. Leedham-Green, P. M. Neumann, and J. Wiegold (1969). This literature assessment is supplied source information, not independently verified here. The displayed strict bound requires an unresolved qualification: a nontrivial abelian group has \(c(P)=1\) and \(b(P)=0\), so the bound as written fails in that case. The supplied record does not specify the necessary restriction or correction.

For \(p=2\), the record attributes to W. Felsch, J. Neub\"user, and W. Plesken (1981) the assertion that for every positive integer \(n\) there exists a finite \(2\)-group \(T_n\) satisfying
\[
c(T_n)\geq b(T_n)+n.
\]
The record writes \(n\in\mathbb{N}\) without specifying whether \(0\) is included; any additional case \(n=0\) follows from the positive cases. These contextual assertions are retained as supplied, and are not assertions of independent literature verification.
\subsection{Short English statement}
The question asks whether every finite group of prime-power order for an odd prime has nilpotency class at most one more than its breadth, where breadth is the base-\(p\) logarithm of the largest conjugacy-class size. The supplied record labels the problem open and reports counterexamples to an analogous bound for \(2\)-groups. Its quoted strict background bound has an unresolved abelian boundary case.
\subsection{Sources}
\begin{itemize}
\item Local UnsolvedMath registry, permanent ID \(701079\), exact source problem code \texttt{catalogue entry}, title \emph{Kourovka Notebook Problem 14.95}; supplied upstream revision \texttt{d371ee917ac9647fbb06cd6fc529deec6f3dda82}. The supplied status is \emph{open}; this is a registry label, not an independently verified literature assessment.
\item \emph{The Kourovka Notebook (Unsolved Problems in Group Theory)}, original problem 14.95, 14th issue (1999); the supplied background also identifies the 21st issue edition (2026), \texttt{arXiv:1401.0300v48}. The registry describes the collection as edited by E. I. Khukhro and V. D. Mazurov. Its proposer field names A. Jaikin-Zapirain, with proposed year 1999; the statement itself credits C. R. Leedham-Green, P. M. Neumann, and J. Wiegold. Both attributions are retained without resolving their respective roles.
\item C. R. Leedham-Green, P. M. Neumann, and J. Wiegold, \emph{J. London Math. Soc. (2)} \textbf{1} (1969), 409--420; cited in the supplied statement for the strict class--breadth bound. No article title is supplied.
\item W. Felsch, J. Neub\"user, and W. Plesken, \emph{J. London Math. Soc. (2)} \textbf{24} (1981), 113--122; cited in the supplied statement for the family of \(2\)-groups. No article title is supplied.
\end{itemize}
\end{document}
